\documentclass[11pt]{article}
\usepackage[utf8]{inputenc}
\usepackage[T1]{fontenc}
\usepackage{amsmath}
\usepackage{amsfonts}
\usepackage{amssymb}
\usepackage{geometry}
\usepackage{hyperref}

\geometry{letterpaper, margin=1in}

\title{\textbf{Problem F: Optimal Mirroring}}
\date{}
\author{}

\begin{document}

\maketitle

\section*{Description}

Nanami Touko is obsessed with the concept of a ``perfect performance.'' To her, identity is like a permutation $a_1, a_2, \dots, a_n$ of $1, 2, \dots, n$. She believes her late sister was the original, perfect permutation $a$. Now, Touko spends her days trying to act in a way that is ``similar'' to that ideal, ensuring that in every scene, her relative ``shape'' matches her sister's perfectly.

Koito Yuu, the only one who truly sees the girl behind the performance, wants to help Touko organize her various scripts. Touko has a collection of $q$ scripts, represented as intervals $[l_i, r_i]$ with $1 \le l_i \le r_i \le n$.

To analyze these performances, Yuu uses the following definitions:

\begin{itemize}
    \item A performance $b_1, b_2, \dots, b_n$ \textbf{holds its shape} for an interval $[l', r']$ if for any pair of integers $(x, y)$ such that $l' \le x < y \le r'$, the relative order of Touko's performance matches the ideal: $b_x < b_y$ if and only if $a_x < a_y$.
    \item A performance $b$ is \textbf{$k$-similar} if it holds its shape for all scripts from $1$ to $k$ (i.e., for all intervals $[l_i, r_i]$ where $1 \le i \le k$).
\end{itemize}

Yuu realizes that there are many possible performances $b$ that could satisfy Touko's strict requirements. To simplify her notes, Yuu decides to encode the set of all $k$-similar performances using Directed Acyclic Graphs (DAG). She defines a \textbf{$k$-encoding} as a DAG $G_k$ on the set of vertices $1, 2, \dots, n$ such that a performance $b$ satisfies $G_k$ if and only if $b$ is $k$-similar.

A performance $b$ \textbf{satisfies} a DAG $G'$ if for every directed edge $u \to v$ in $G'$, we have $b_u < b_v$.

Since Yuu is the only one who can truly understand Touko, she wants to figure out the minimum number of edges among all $k$-encodings for each $k$ from $1$ to $q$.

\section*{Input}

The first line contains two integers $n$ and $q$ ($1 \le n \le 25000$, $1 \le q \le 100000$).

The second line contains $n$ integers $a_1, a_2, \dots, a_n$ which form a permutation of $1, 2, \dots, n$ (Touko's sister's ideal performance).

The next $q$ lines each contain two integers $l_i$ and $r_i$ ($1 \le l_i \le r_i \le n$), representing the scripts.

\section*{Output}

Print $q$ lines. The $k$-th line should contain a single integer --- the minimum number of edges in an optimal $k$-encoding.

\section*{Examples}

\subsection*{Example 1}
\textbf{Input:}
\begin{verbatim}
4 3
2 4 1 3
1 3
2 4
1 4
\end{verbatim}

\textbf{Output:}
\begin{verbatim}
2
4
3
\end{verbatim}

\section*{Note}

In the first test case:
\begin{itemize}
    \item For $k=1$, the performance must mirror the interval $[1,3]$. The ideal order is $a_1=2, a_2=4, a_3=1$, meaning $a_3 < a_1 < a_2$. The optimal encoding $G_1$ has edges $3 \to 1$ and $1 \to 2$. Total 2 edges.
    \item For $k=2$, Touko must mirror $[1,3]$ and $[2,4]$.
    \item For $k=3$, adding $[1,4]$ might actually make some previous edges redundant (transitive reduction).
\end{itemize}

\end{document}